\textbf{Proof.}
The published comparator supplies only the variance-ratio normalization and maximin operation already fixed in \(\mathrm{def:ryan\text{-}relative\text{-}efficiency}\).  The calculation below concerns this paper's finite menu.  By \(\mathrm{thm:quadruple\text{-}menu\text{-}frontier}\), \(r_j^\star=4c_j/n^2\).  For any feasible design, its trace certificates imply, with \(p_k=p_k(\pi)\),
\[
\operatorname{RE}_1(\pi)\le\frac1{1+F_1},
\qquad
\operatorname{RE}_2(\pi)\le\frac1{1+F_2},
\tag{1}
\]
where
\[
F_1=\kappa_1(ap_2+bp_3),
\qquad
F_2=\kappa_2(ap_1+bp_3).
\tag{2}
\]
All denominators are positive because \(a,b,\kappa_j>0\).  It therefore suffices to minimize \(\max(F_1,F_2)\) and then exhibit a law attaining both trace bounds.

Assume \(a>2b>0\) and \(\kappa_2>\kappa_1>0\), and define
\[
D=\kappa_1(a-b)+\kappa_2b,
\qquad
\alpha=\frac{\kappa_2b}{D},
\qquad
1-\alpha=\frac{\kappa_1(a-b)}D.
\tag{3}
\]
Since \(a-b>b>0\), \(D>0\) and \(0<\alpha<1\).  Substituting \(p_3=1-p_1-p_2\) into (2) gives the affine lower certificate
\[
\alpha F_1+(1-\alpha)F_2
=\frac{ab\kappa_1\kappa_2}{D}
+\frac{a(a-2b)\kappa_1\kappa_2}{D}p_1
\ge \frac{ab\kappa_1\kappa_2}{D}=:m.
\tag{4}
\]
Because a maximum dominates every convex combination, \(\max(F_1,F_2)\ge m\).

Set
\[
p_1^\circ=0,
\qquad
p_2^\circ=\frac{b(\kappa_2-\kappa_1)}D,
\qquad
p_3^\circ=1-p_2^\circ.
\tag{5}
\]
The numerator of \(p_2^\circ\) is positive and
\[
D-b(\kappa_2-\kappa_1)=a\kappa_1>0,
\tag{6}
\]
so (5) lies in the simplex.  Direct substitution gives \(F_1=F_2=m\).  Choose directions independently across blocks with probabilities (5), together with independent symmetric signs.  The law is in \(\Pi_n\), and its independent centered block scores give
\[
C_j=4c_j(1+F_j)I_q,
\qquad
R_j=\frac{4c_j}{n^2}(1+F_j).
\tag{7}
\]
Thus both inequalities in (1) are equalities, and
\[
\eta_{\mathrm{REL}}
=\frac1{1+m}
=\left[1+\frac{ab\kappa_1\kappa_2}
{\kappa_1(a-b)+\kappa_2b}\right]^{-1}.
\tag{8}
\]

If an arbitrary design attains (8), equality is necessary in both
\[
\min_j\operatorname{RE}_j(\pi)
\le\{1+\max(F_1,F_2)\}^{-1}
\le(1+m)^{-1}.
\tag{9}
\]
The coefficient of \(p_1\) in (4) is strictly positive, so equality forces \(p_1=0\).  Since \(0<\alpha<1\), equality between the maximum and the convex combination forces \(F_1=F_2=m\).  Solving on the face \(p_1=0\) gives exactly (5).

Finally, the \(a>2b\) branch of the gain-normalized frontier says every gain-normalized optimizer has \(p_3=1\), whereas (5) has \(p_2^\circ>0\).  The optimizer sets are therefore disjoint.  The comparison uses the same fixed class risks, oracle risks, and blinded assignment space; only the normalization differs, so no substantive submodel claim about the comparator is needed.  \(\square\)